\documentclass[11pt]{article}

\usepackage[margin=1in]{geometry}
\usepackage{graphicx}
\usepackage{subcaption}
\usepackage{booktabs}
\usepackage{hyperref}
\usepackage{algorithm}
\usepackage{algorithmic}

\usepackage{amsmath,amssymb,mathtools,amsthm}
\usepackage[capitalize,noabbrev]{cleveref}
\crefname{assumption}{Assumption}{Assumptions}
\usepackage[numbers,sort&compress]{natbib}

\hypersetup{
  colorlinks=true,
  linkcolor=blue,
  citecolor=blue,
  urlcolor=blue
}

%%%%%%%%%%%%%%%%%%%%%%%%%%%%%%%%
% THEOREMS (OR-style: section-numbered)
%%%%%%%%%%%%%%%%%%%%%%%%%%%%%%%%
\theoremstyle{plain}
\newtheorem{theorem}{Theorem}[section]
\newtheorem{proposition}[theorem]{Proposition}
\newtheorem{lemma}[theorem]{Lemma}
\newtheorem{corollary}[theorem]{Corollary}
\theoremstyle{definition}
\newtheorem{definition}[theorem]{Definition}
\newtheorem{assumption}[theorem]{Assumption}
\newtheorem{example}[theorem]{Example}
\theoremstyle{remark}
\newtheorem{remark}[theorem]{Remark}

\title{Leveraging Inference-Time Compute for Diffusion Models via Global Scheduling of Denoising Trajectories}
\author{Anonymous Author(s)}
\date{}

% === Theory Extract: auto-generated, do not edit ===
\usepackage{xcolor}
\definecolor{srcgray}{gray}{0.5}
\newcommand{\srcfile}[1]{\hfill{\tiny\color{srcgray}[\detokenize{#1}]}}

\title{Theory Extract: GAINS}
\author{Auto-generated from main paper}
\date{\today}

\begin{document}
\maketitle

\tableofcontents
\newpage

% ============================================
% Stub labels for cross-references to main paper
% (equations, figures, algorithms, sections)
% ============================================
\section*{Cross-Reference Stubs}
{\small The following items are referenced by the
theory but defined in the main paper.}
\medskip

% Stub algorithm from algorithm.tex
\begin{algorithm}[h]\caption{See main paper.}\label{alg:offline_online_budget}\end{algorithm}
% Stub equation from algorithm.tex
\begin{equation} \text{(see main paper)} \label{eq:alloc} \end{equation}
% Stub equation from preliminaries.tex
\begin{equation} \text{(see main paper)} \label{eq:ddpm-reverse} \end{equation}
% Stub equation from general_local_search.tex
\begin{equation} \text{(see main paper)} \label{eq:gain-general} \end{equation}
% Stub equation from preliminaries.tex
\begin{equation} \text{(see main paper)} \label{eq:sgm-langevin} \end{equation}
% Stub figure from algorithm.tex
\begin{figure}[h]\centering\caption{See main paper.}\label{fig:stopping}\end{figure}
% Stub section from framework.tex
\subsection{(See main paper)} \label{sec:framework}
% Stub section from appendix_mdp.tex
\subsection{(See main paper)} \label{sec:mdp}
% Stub section from appendix_proofs.tex
\subsection{(See main paper)} \label{sec:proof-offline}
% Stub section from appendix_proofs.tex
\subsection{(See main paper)} \label{sec:proof-online}
\newpage

\section{Preliminaries}
\srcfile{preliminaries.tex}

\begin{remark}
Both DDPM and SGM traverse $T$ discrete timesteps.
At each step, the transition from $\mathbf{x}_t$ to $\mathbf{x}_{t-1}$ depends
on a stochastic noise variable (the injected $\boldsymbol{\varepsilon}$ in
\eqref{eq:ddpm-reverse} or the Langevin noise in~\eqref{eq:sgm-langevin}).
This shared structure underpins the transition abstraction
in \cref{sec:framework}.
\end{remark}

\begin{definition}[Noise Trajectory Search Problem]
\label{def:nts}
Given a pretrained diffusion model with transition function
$F_\theta$, a conditioning signal $c$, a verifier $v(\cdot, c)$, and a total
compute budget of $B$ function evaluations (NFE), the \emph{noise trajectory
search problem} is to find a sequence of noise variables
$\{\varepsilon_t^*\}_{t=1}^T$ that maximizes the verifier score of the final
generated sample:
\begin{equation}
  \max_{\{\varepsilon_t\}_{t=1}^T}\;
  v\!\bigl(x_0(\varepsilon_1, \ldots, \varepsilon_T),\, c\bigr)
  \quad \text{subject to} \quad
  \text{total NFE} \leq B.
  \label{eq:nts-problem}
\end{equation}
\end{definition}

\begin{definition}[Verifier]
\label{def:verifier}
A \emph{verifier} $v : \mathcal{X} \times \mathcal{C} \to \mathbb{R}$ scores a
generated sample $x_0$ against its conditioning signal $c$, returning a scalar
quality measure.
The verifier is fixed and not trained during search.
Examples include perceptual quality metrics, task-specific reward functions,
and domain-specific performance measures.
\end{definition}

\begin{definition}[Function Evaluation Budget]
\label{def:nfe}
A single \emph{function evaluation} (NFE) corresponds to one forward pass
through the denoising network.
The budget $B$ caps how many forward passes the sampler may
perform when generating a single sample.
Standard diffusion sampling consumes exactly $T$ NFE (one per timestep);
noise trajectory search spends the remaining $B - T$ NFE on local noise
refinement.
\end{definition}

\section{GAINS Algorithm and Theory}
\srcfile{algorithm.tex}

\begin{assumption}[Verifier smoothness]
\label{asm:smooth}
The composed verifier score function
$\Psi_t(\cdot\,; c) \in C^2(\mathbb{R}^d)$, and there exists a constant
$L_t < \infty$ such that
$\|\nabla^2 \Psi_t(x; c)\|_{\mathrm{op}} \le L_t$
for all $x$ in a neighborhood of the line segment between $\bar{X}_{t-h}$
and $X_{t-h}$, almost surely.
\end{assumption}

\begin{proposition}[Score Taylor expansion]
\label{prop:taylor}
Under \cref{asm:smooth}, conditional on $(X_t, c)$,
\begin{equation}
  S_t \;=\; \Psi_t(\bar{X}_{t-h};\, c)
        \;+\; g_t\sqrt{h}\,
              \bigl\langle \nabla\Psi_t(\bar{X}_{t-h};\, c),\, \xi \bigr\rangle
        \;+\; R_t,
  \label{eq:taylor-score}
\end{equation}
where the remainder satisfies
$|R_t| \le \tfrac{1}{2}\,L_t\,g_t^2\,h\,\|\xi\|^2$ almost surely,
and hence
$\mathbb{E}[\,|R_t|\mid X_t, c\,] \le \tfrac{1}{2}\,L_t\,g_t^2\,h\,d$.
\end{proposition}

\begin{proof}
By construction,
$X_{t-h} = \bar{X}_{t-h} + g_t\sqrt{h}\,\xi$.
A second-order Taylor expansion of $\Psi_t(\cdot\,; c)$ around
$\bar{X}_{t-h}$ gives
\[
  \Psi_t(X_{t-h}; c)
  = \Psi_t(\bar{X}_{t-h}; c)
    + \bigl\langle \nabla\Psi_t(\bar{X}_{t-h}; c),\,
        X_{t-h} - \bar{X}_{t-h} \bigr\rangle
    + R_t,
\]
where
$R_t = \tfrac{1}{2}(X_{t-h} - \bar{X}_{t-h})^\top
  \nabla^2\Psi_t(\widetilde{X}; c)\,
  (X_{t-h} - \bar{X}_{t-h})$
for some $\widetilde{X}$ on the segment joining $\bar{X}_{t-h}$ and $X_{t-h}$.
Substituting $X_{t-h} - \bar{X}_{t-h} = g_t\sqrt{h}\,\xi$ yields
\eqref{eq:taylor-score}.
The remainder bound follows from
$|R_t| \le \tfrac{1}{2}\|\nabla^2\Psi_t(\widetilde{X}; c)\|_{\mathrm{op}}
  \,\|g_t\sqrt{h}\,\xi\|^2
  \le \tfrac{1}{2}L_t\,g_t^2\,h\,\|\xi\|^2$.
Taking conditional expectations and using
$\mathbb{E}[\|\xi\|^2] = d$ completes the proof.
\end{proof}

\begin{theorem}[Location-scale structure of per-step scores]
\label{thm:loc-scale}
Define
\begin{equation}
  \mu_t \;:=\; \Psi_t(\bar{X}_{t-h};\, c),
  \qquad
  \sigma_t \;:=\; g_t\sqrt{h}\,\|\nabla\Psi_t(\bar{X}_{t-h};\, c)\|.
  \label{eq:mu-sigma-sde}
\end{equation}
Then, conditional on $(X_t, c)$:
\begin{enumerate}
\item[\textbf{(i)}]
  There exists $Z_t \sim \mathcal{N}(0,1)$ such that
  $S_t = \mu_t + \sigma_t\,Z_t + R_t$.
\item[\textbf{(ii)}]
  As $g_t\sqrt{h} \to 0$,
  \begin{equation}
    \operatorname{Var}(S_t \mid X_t, c)
    \;=\; \sigma_t^2 + O_d(g_t^3\,h^{3/2}).
    \label{eq:var-approx}
  \end{equation}
\item[\textbf{(iii)}]
  Let $S_t^{(1)}, \ldots, S_t^{(K)}$ be $K$ i.i.d.\ copies of $S_t$
  generated by i.i.d.\ noise draws $\xi_1, \ldots, \xi_K$.
  Then
  \begin{equation}
    G_t(K)
    \;:=\; \mathbb{E}\Bigl[\max_{1 \le k \le K} S_t^{(k)}
           \,\Big|\, X_t, c\Bigr] - \mu_t
    \;=\; \sigma_t\,a_K + O_{K,d,L_t}(g_t^2\,h),
    \label{eq:gain-sde}
  \end{equation}
  where $a_K := \mathbb{E}[\max(Z_1, \ldots, Z_K)]$ with
  $Z_k \stackrel{\mathrm{iid}}{\sim} \mathcal{N}(0,1)$.
  In particular,
  \begin{equation}
    G_t(K) \;=\; g_t\sqrt{h}\,\|\nabla\Psi_t(\bar{X}_{t-h};\, c)\|\,a_K
              \;+\; O_{K,d,L_t}(g_t^2\,h).
    \label{eq:gain-explicit}
  \end{equation}
\end{enumerate}
\end{theorem}

\begin{proof}
For part~(i), when $\|\nabla\Psi_t(\bar{X}_{t-h}; c)\| > 0$, write
$\langle \nabla\Psi_t, \xi \rangle
= \|\nabla\Psi_t\|\,Z_t$
where $Z_t := \langle \nabla\Psi_t / \|\nabla\Psi_t\|, \xi \rangle$.
Since $\nabla\Psi_t / \|\nabla\Psi_t\|$ is a fixed unit vector conditional on
$(X_t, c)$ and $\xi \sim \mathcal{N}(0, I)$, the projection $Z_t$ is standard
normal: $Z_t \sim \mathcal{N}(0,1)$.
(When $\nabla\Psi_t(\bar{X}_{t-h}; c) = 0$, we have $\sigma_t = 0$ and the
linear term vanishes; the result holds trivially with $Z_t$ defined
arbitrarily.)

For part~(ii), expand
$\operatorname{Var}(S_t \mid X_t, c)
= \operatorname{Var}(\sigma_t Z_t + R_t)
= \sigma_t^2 + 2\,\operatorname{Cov}(\sigma_t Z_t, R_t) + \operatorname{Var}(R_t)$.
By Cauchy--Schwarz,
\[
  |\operatorname{Cov}(\sigma_t Z_t, R_t)|
  \le \sqrt{\operatorname{Var}(\sigma_t Z_t)}\,
      \sqrt{\operatorname{Var}(R_t)}
  \le \sigma_t\,\sqrt{\mathbb{E}[R_t^2]}.
\]
Using \cref{prop:taylor},
$|R_t| \le \tfrac{1}{2}L_t g_t^2 h \|\xi\|^2$, hence
\[
  \mathbb{E}[R_t^2 \mid X_t, c]
  \le \tfrac{1}{4}L_t^2 g_t^4 h^2\,\mathbb{E}[\|\xi\|^4]
  = \tfrac{1}{4}L_t^2 g_t^4 h^2\,d(d+2),
\]
because $\xi\sim\mathcal{N}(0,I)$ implies
$\mathbb{E}[\|\xi\|^4]=d(d+2)$.
Therefore
\[
  |\operatorname{Cov}(\sigma_t Z_t, R_t)|
  \le \tfrac{1}{2}L_t\sqrt{d(d+2)}\,\sigma_t g_t^2 h
  = O_d(g_t^3 h^{3/2}),
\]
with the implicit constant depending on the conditional local geometry
through $L_t$ and $\|\nabla\Psi_t\|$.
Likewise,
\[
  \operatorname{Var}(R_t)
  \le \mathbb{E}[R_t^2]
  \le \tfrac{1}{4}L_t^2 g_t^4 h^2\,d(d+2)
  = O_d(g_t^4 h^2),
\]
which is lower order than the covariance term as $g_t\sqrt{h}\to 0$.

For part (iii), write
$S_t^{(k)} = \mu_t + \sigma_t Z_k + R_t^{(k)}$
for each candidate $k$.
Then
$\max_k S_t^{(k)} = \mu_t + \max_k (\sigma_t Z_k + R_t^{(k)})$.
Using the elementary inequality
$|\max_k (a_k + b_k) - \max_k a_k| \le \max_k |b_k|$,
we obtain
$|\max_k S_t^{(k)} - (\mu_t + \sigma_t \max_k Z_k)|
\le \max_k |R_t^{(k)}|$.
Since $\max_k |R_t^{(k)}| \le \sum_k |R_t^{(k)}|$,
taking conditional expectations gives
$\mathbb{E}[\max_k |R_t^{(k)}| \mid X_t, c]
\le K \cdot \tfrac{1}{2} L_t g_t^2 h\,d$.
Therefore
$G_t(K) = \sigma_t\,a_K + O_{K,d,L_t}(g_t^2 h)$,
where the implicit constant in $O_{K,d,L_t}(\cdot)$ depends on $K$, the
dimension~$d$, and the Hessian bound~$L_t$ from \cref{asm:smooth}, but not
on the diffusion coefficient $g_t$ or the verifier gradient
$\nabla\Psi_t$.
Substituting $\sigma_t = g_t\sqrt{h}\,\|\nabla\Psi_t(\bar{X}_{t-h}; c)\|$
yields~\eqref{eq:gain-explicit}.
\end{proof}

\begin{remark}
\cref{thm:loc-scale} identifies the precise quantities governing timestep
sensitivity:
$\sigma_t = g_t\sqrt{h}\,\|\nabla\Psi_t(\bar{X}_{t-h}; c)\|$
multiplies the diffusion coefficient $g_t$, the step size
$\sqrt{h}$, and the verifier gradient norm at the
deterministic Euler point.
Timesteps with large $g_t$ (strong noise injection) or large
$\|\nabla\Psi_t\|$ (steep verifier landscape) benefit most from noise
search, explaining the non-uniform sensitivity patterns observed in offline profiling.
\end{remark}

\begin{corollary}[SDE-based approximate allocation]
\label{cor:sde-alloc}
Under \cref{asm:smooth}, with step size $h$ and diffusion
coefficients $\{g_t\}_{t=1}^T$, for any fixed budget allocation
$\{K_t\}_{t=1}^T$,
\begin{equation}
  \sum_{t=1}^{T} G_t(K_t)
  \;=\; \sum_{t=1}^{T} g_t\sqrt{h}\,
        \|\nabla\Psi_t(\bar{X}_{t-h}; c)\|\,a_{K_t}
  \;+\; \sum_{t=1}^{T} O_{K_t,d,L_t}(g_t^2\,h).
  \label{eq:aggregate-gain}
\end{equation}
Dropping the second-order remainder reduces the offline allocation problem
to
\begin{equation}
  \max_{\{K_t\}_{t=1}^T:\;\sum_t K_t = B}\;
  \sum_{t=1}^{T} g_t\sqrt{h}\,
  \|\nabla\Psi_t(\bar{X}_{t-h}; c)\|\,a_{K_t}.
  \label{eq:alloc-sde}
\end{equation}
\end{corollary}

\begin{proposition}[Offline optimality under independence]
\label{prop:offline}
Suppose $\sigma_t$ does not depend on the prompt or search history
(i.e., $\sigma_t = \sigma_t$ for all $c, x_t$).
Then:
\begin{enumerate}
\item[\textbf{(i)}]
  Problem~\eqref{eq:alloc} decomposes across prompts: any optimal
  allocation $\{K_t^*\}$ is simultaneously optimal for every instance.
\item[\textbf{(ii)}]
  The optimal allocation equalizes marginal gains:
  $\sigma_t\,a_{K_t^*}' = \lambda^*$ for all~$t$ with $K_t^* > 1$,
  where $\lambda^*$ is the shadow price of the budget constraint.
  Consequently, $K_t^*$ is increasing in~$\sigma_t$ (water-filling).
\item[\textbf{(iii)}]
  $\sum_t \sigma_t\,a_{K_t^*} > \sum_t \sigma_t\,a_{B/T}$ whenever
  $\{\sigma_t\}$ are not all equal, and the gap grows with the
  dispersion of~$\{\sigma_t\}$.
\end{enumerate}
\end{proposition}

\begin{remark}
\cref{prop:offline} justifies offline profiling.
Part~(i) guarantees that under the independence condition, a single fixed
allocation optimizes every prompt simultaneously, so profiling on a small calibration
set generalizes.
Part~(ii) reveals \emph{water-filling} structure: timesteps with larger
$\sigma_t = g_t\sqrt{h}\,\|\nabla\Psi_t\|$ (where noise perturbations
shift the verifier score most) receive proportionally more budget.
Part~(iii) quantifies the non-uniform advantage: the more
heterogeneous the per-step sensitivities, the larger the gain
over uniform scheduling.
\end{remark}

\begin{proposition}[Online adaptation under prompt dependence]
\label{prop:online}
Suppose $\sigma_t = \sigma_t(c, x_t)$ varies with the prompt and search
history.
Define the oracle value
$V^*(\boldsymbol{\sigma})=\max_{\{K_t\}:\sum_t K_t=B}
\sum_t\sigma_t\,a_{K_t}$
for variance profile $\boldsymbol{\sigma}=(\sigma_1,\dots,\sigma_T)$.
Then:
\begin{enumerate}
\item[\textbf{(i)}]
  $V^*$ is convex in $\boldsymbol{\sigma}$.
  The expected regret of any fixed allocation satisfies
  \begin{equation}
    \mathbb{E}_{\boldsymbol{\sigma}}\!\bigl[V^*(\boldsymbol{\sigma})\bigr]
    \;-\; V^*\!\bigl(\bar{\boldsymbol{\sigma}}\bigr) \;\ge\; 0,
    \label{eq:jensen-gap}
  \end{equation}
  where $\bar{\boldsymbol{\sigma}}=\mathbb{E}[\boldsymbol{\sigma}]$,
  with equality whenever the same allocation is optimal almost surely on
  the support of $\boldsymbol{\sigma}$.
\item[\textbf{(ii)}]
  For integer budgets, the one-step marginal gain obeys
  \begin{equation}
    \Delta_t(K)
    \;:=\;
    G_t(K{+}1)-G_t(K)
    \;=\;
    \underbrace{\sigma_t}_{\text{sensitivity}}
    \;\cdot\; \underbrace{\Delta a_K}_{\text{exhaustion}}
    \;+\; O_{K,d,L_t}(g_t^2 h),
    \qquad
    \Delta a_K:=a_{K+1}-a_K .
    \label{eq:marginal-decomp}
  \end{equation}
  The optimal stopping condition is $\Delta_t(K)\le\lambda^*$.
  The running mean $\bar{g}$ upper-bounds the active discrete marginals
  up to the same $O_{K,d,L_t}(g_t^2 h)$ remainder, so
  $\beta_g\bar{g}$ serves as an adaptive estimate of~$\lambda^*$.
  The gain condition $\widetilde{g}_t^{(j)}<\beta_g\bar{g}$ tests the
  leading-order product $\sigma_t\,\Delta a_K$ against the
  shadow-price scale~$\lambda^*$; the
  variance condition is an empirical proxy for low sensitivity rather
  than an exact test of $\sigma_t$.
  Requiring both conditions prevents premature stopping at
  high-value timesteps.
\end{enumerate}
\end{proposition}

\begin{remark}
\cref{prop:online} motivates the online controller.
Part~(i) shows that $V^*$, as a pointwise maximum of linear functions, is convex in the variance
profile, so by Jensen's
inequality any fixed allocation tuned to the population mean
$\bar{\boldsymbol{\sigma}}$ incurs a gap
(\cref{eq:jensen-gap}) growing with cross-instance variability
of~$\boldsymbol{\sigma}$.
Equality holds whenever the same allocation remains optimal across the
support of $\boldsymbol{\sigma}$; a constant profile is only a
sufficient special case.
Part~(ii) decomposes the marginal gain~\eqref{eq:marginal-decomp}
into two factors, each monitored by a separate
threshold in the online controller.
Low gain with high variance ($\Delta a_K$ small, $\sigma_t$ large) signals
a responsive but exhausted timestep where a lucky draw may still yield a large jump, so stopping is
premature.
Moderate gain with low variance ($\Delta a_K$ moderate, $\sigma_t$ small)
signals inherent insensitivity where gains plateau
regardless, yet search remains productive at this iteration count.
The variance statistic used by the controller should therefore be read
as a proxy for the latent sensitivity $\sigma_t$, justified by the
location-scale variance approximation in \cref{thm:loc-scale}(ii),
rather than as an exact measurement of~$\sigma_t$ itself.
Only when \emph{both} factors are small is stopping unambiguously
justified, as the two-threshold diagram in \cref{fig:stopping} shows.
\end{remark}

\section{General Local Search Operators}
\srcfile{general_local_search.tex}

\begin{assumption}[Local search regularity]
\label{asm:local-search}
A local search operator $\mathcal{L}$ satisfies:
\begin{enumerate}
\item[\textbf{(A1)}]
  \emph{Strict monotone improvement.}\;
  $G_t^{\mathcal{L}}(K)$ is strictly increasing in~$K$ for $K\ge 1$.
\item[\textbf{(A2)}]
  \emph{Diminishing returns.}\;
  $G_t^{\mathcal{L}}(K{+}1)-G_t^{\mathcal{L}}(K)$ is bounded above
  by a non-increasing sequence, i.e., the marginal gains do not
  accelerate.
\item[\textbf{(A3)}]
  \emph{Sensitivity scaling.}\;
  In the linearized regime ($g_t\sqrt{h}\to 0$),
  \begin{equation}
    G_t^{\mathcal{L}}(K)
    \;=\; \sigma_t\,\phi_K^{\mathcal{L}}
    \;+\; O_{K,d,L_t}(g_t^2\,h),
    \label{eq:sensitivity-scaling}
  \end{equation}
  where $\phi_K^{\mathcal{L}}\ge 0$ depends on the operator
  $\mathcal{L}$ and the budget $K$ but \emph{not} on the
  timestep~$t$, and the implicit constant in $O_{K,d,L_t}(\cdot)$ may
  depend on $K$, $d$, and $L_t$.
\item[\textbf{(A4)}]
  \emph{Rotational equivariance.}\;
  For any orthogonal matrix $Q\in\mathbb{R}^{d\times d}$, the joint
  distribution of $\{Q\xi^{(j)}\}_{j\ge 0}$ under $\mathcal{L}$
  equals that of the iterates produced by $\mathcal{L}$ starting
  from $Q\xi^{(0)}$.
\end{enumerate}
\end{assumption}

\begin{remark}[Grounding of assumptions]
\label{rmk:grounding}
Each assumption captures a natural property of practical operators.
\textbf{(A1)} holds whenever the operator
maintains a running best-so-far (a ``keep best'' wrapper) and each
iteration has a positive probability of improving the incumbent.
\textbf{(A2)} reflects the general principle that the first
few iterations are most productive: for random search this follows from
order-statistic theory ($a_K$ is concave;
\citealt{david2003order}, Ch.~4); for gradient-based methods it
follows from convergence on smooth objectives; for MCMC it follows from
mixing-time arguments.
\textbf{(A3)} arises because, in the small-noise regime,
$\Psi_t(\bar{X}_{t-h}+g_t\sqrt{h}\,\xi;c)
= \mu_t+\sigma_t\langle u_t,\xi\rangle + O(g_t^2 h)$
where $u_t:=\nabla\Psi_t/\|\nabla\Psi_t\|$ is a fixed unit vector
(\cref{prop:taylor}).
Any local search on $\xi$ therefore reduces to optimizing a linear
function $f(\xi)=\langle u,\xi\rangle$, and the gain from this
optimization depends on the search method and budget but not on the
direction~$u$ or the scaling factor~$\sigma_t$, beyond a
multiplicative factor.
\textbf{(A4)} holds for all three operators below, since each uses
only isotropic primitives (Gaussian exploration draws and
uniform-direction local perturbations); it ensures that the
projected gain $\mathbb{E}[\max_j\langle u,\xi^{(j)}\rangle]$ does
not depend on the direction~$u$.
\end{remark}

\begin{theorem}[General gain factorization]
\label{thm:general-gain}
Under \cref{asm:smooth,asm:local-search}, for any local search operator
$\mathcal{L}$ with operator-specific gain sequence
$\{\phi_K^{\mathcal{L}}\}_{K\ge 0}$, the expected gain from $K$
iterations at timestep~$t$ satisfies
\begin{equation}
  G_t^{\mathcal{L}}(K)
  \;=\;
  \sigma_t\,\phi_K^{\mathcal{L}}
  + O_{K,d,L_t}(g_t^2\,h),
  \label{eq:general-gain-factorization}
\end{equation}
where $\sigma_t=g_t\sqrt{h}\,\|\nabla\Psi_t(\bar{X}_{t-h};c)\|$
and the implicit constant in $O_{K,d,L_t}(\cdot)$ depends on $K$, $d$,
and $L_t$ but not on $g_t$ or $\nabla\Psi_t$.
The sequence $\phi_K^{\mathcal{L}}$ is strictly increasing for
$K\ge 1$ (by~A1), has non-accelerating marginal gains (by~A2),
and satisfies $\phi_0^{\mathcal{L}}=0$.
\end{theorem}

\begin{proof}
By \cref{prop:taylor},
$\Psi_t(\bar{X}_{t-h}+g_t\sqrt{h}\,\xi;c)
=\mu_t+\sigma_t\langle u_t,\xi\rangle+R_t(\xi)$,
where $u_t=\nabla\Psi_t(\bar{X}_{t-h};c)/\|\nabla\Psi_t(\bar{X}_{t-h};c)\|$
and $|R_t(\xi)|\le\tfrac{1}{2}L_t g_t^2 h\|\xi\|^2$.
Since $\mu_t$ is constant across iterates, substituting
into~\eqref{eq:gain-general} gives
\begin{align*}
  G_t^{\mathcal{L}}(K)
  &= \mathbb{E}\Bigl[
       \max_{0\le j\le K}
       \bigl(\sigma_t\langle u_t,\xi^{(j)}\rangle + R_t(\xi^{(j)})\bigr)
     \,\Big|\,X_t,c\Bigr] \\
  &= \sigma_t\,\mathbb{E}\Bigl[
       \max_{0\le j\le K}\langle u_t,\xi^{(j)}\rangle
     \,\Big|\,X_t,c\Bigr]
     + O_{K,d,L_t}(g_t^2 h),
\end{align*}
where the second equality uses the elementary inequality
$|\max_j(a_j+b_j)-\max_j a_j|\le\max_j|b_j|$
(as in \cref{thm:loc-scale}(iii)), the non-negativity $\sigma_t\ge 0$
to factor $\sigma_t$ out of the max, and the remainder bound
$\mathbb{E}[\max_{0\le j\le K}|R_t(\xi^{(j)})|]
\le\tfrac{1}{2}L_t g_t^2 h\,\mathbb{E}[\max_{0\le j\le K}\|\xi^{(j)}\|^2]$.
The subscript in $O_{K,d,L_t}(\cdot)$ reflects that the constant depends on
$\mathbb{E}[\max_j\|\xi^{(j)}\|^2]$, which may grow with~$K$ for
operators whose iterates accumulate drift.

By~(A4), the joint distribution of $\{\xi^{(j)}\}$ is rotationally
equivariant.
Hence, for any fixed unit vector $u_t$, the projected iterates
$\{\langle u_t,\xi^{(j)}\rangle\}$ have the same distribution as
$\{\langle e_1,Q\xi^{(j)}\rangle\}$ for any orthogonal $Q$ satisfying
$Qu_t=e_1$.
By~(A4), the rotated process $\{Q\xi^{(j)}\}$ has the same law as the
operator started from $Q\xi^{(0)}$; because
$\xi^{(0)}\sim\mathcal{N}(0,I)$, the initialization is itself
rotationally invariant, so $Q\xi^{(0)}\stackrel{d}{=}\xi^{(0)}$.
Therefore the full iterate sequence has the same law after rotation, and
hence $\{\langle u_t,\xi^{(j)}\rangle\}$ has the same distribution as
$\{\langle e_1,\xi^{(j)}\rangle\}$ (where $e_1$ is any fixed
coordinate direction).
Therefore the expected projected gain
$\mathbb{E}[\max_{0\le j\le K}\langle u_t,\xi^{(j)}\rangle]$
depends on $\mathcal{L}$ and $K$ only; we denote it
$\phi_K^{\mathcal{L}}$.
Strict monotonicity of $\phi_K^{\mathcal{L}}$
follows from~(A1), non-accelerating marginals from~(A2),
and $\phi_0^{\mathcal{L}}=0$ since
$\mathbb{E}[\langle u,\xi^{(0)}\rangle]=0$ for
$\xi^{(0)}\sim\mathcal{N}(0,I)$.
\end{proof}

\begin{corollary}[Offline water-filling for general operators]
\label{cor:offline-general}
Under the conditions of \cref{thm:general-gain}, with
prompt-independent~$\sigma_t$,
the allocation problem
\begin{equation}
  \max_{\{K_t\}:\,\sum_t K_t=B,\;K_t\ge 1}\;
  \sum_{t=1}^{T}\sigma_t\,\phi_{K_t}^{\mathcal{L}}
  \label{eq:alloc-general}
\end{equation}
has the same structure as \cref{prop:offline}:
\begin{enumerate}
\item[\textup{(i)}]
  The optimal allocation $\{K_t^*\}$ is simultaneously optimal for
  every instance.
\item[\textup{(ii)}]
  There exists an optimal allocation that is non-decreasing
  in~$\sigma_t$.
  When $\phi_K$ is strictly concave, the active timesteps satisfy
  the strict water-filling ordering.
  When $\phi_K$ is linear, every optimizer concentrates the
  allocable budget on $\arg\max_t\sigma_t$.
\item[\textup{(iii)}]
  Non-uniform allocation weakly dominates uniform; the
  dominance is strict whenever $\{\sigma_t\}$ are not all equal
  and $\phi_K$ is strictly concave.
\end{enumerate}
\end{corollary}

\begin{proof}
The proof splits by curvature.
Because \cref{thm:general-gain} gives a common, $t$-independent gain
sequence $\phi_K^{\mathcal{L}}$ with non-accelerating marginals, the
discrete exchange argument from \cref{prop:offline} still shows that an
optimal allocation can be chosen monotone in~$\sigma_t$ and that moving
budget from a lower-$\sigma_t$ timestep to a higher-$\sigma_t$ timestep
never hurts, establishing parts~(i) and~(iii) at the weak level.
When $\phi_K$ is strictly concave, the KKT argument in
\cref{sec:proof-offline} applies verbatim with $\phi_K$ in place of $a_K$,
yielding the strict water-filling ordering among active timesteps.
When $\phi_K$ is linear, say $\phi_K=\alpha K$ with $\alpha>0$,
\eqref{eq:alloc-general} becomes a linear program
$\alpha\sum_t \sigma_t K_t$, so every optimizer assigns all budget beyond
the mandatory floor $K_t\ge 1$ to $\arg\max_t \sigma_t$.
\end{proof}

\begin{corollary}[Online Jensen gap and dual-threshold stopping for general operators]
\label{cor:online-general}
Under the conditions of \cref{thm:general-gain}, with
prompt-dependent $\sigma_t=\sigma_t(c,x_t)$,
define the generalized oracle
\begin{equation}
  V_{\mathcal{L}}^*(\boldsymbol{\sigma})
  \;=\;
  \max_{\{K_t\}:\,\sum_t K_t=B,\;K_t\ge 1}\;
  \sum_{t=1}^{T}\sigma_t\,\phi_{K_t}^{\mathcal{L}}.
  \label{eq:oracle-general}
\end{equation}
Then:
\begin{enumerate}
\item[\textup{(i)}]
  $V_{\mathcal{L}}^*$ is convex in $\boldsymbol{\sigma}$, and
  \begin{equation}
    \mathbb{E}_{\boldsymbol{\sigma}}\!
    \bigl[V_{\mathcal{L}}^*(\boldsymbol{\sigma})\bigr]
    \;\ge\;
    V_{\mathcal{L}}^*\!\bigl(\bar{\boldsymbol{\sigma}}\bigr),
    \label{eq:jensen-general}
  \end{equation}
  with equality whenever the same allocation is optimal almost surely on
  the support of $\boldsymbol{\sigma}$.
\item[\textup{(ii)}]
  For integer budgets, the one-step marginal gain satisfies
  \begin{equation}
    \Delta_t^{\mathcal{L}}(K)
    \;:=\;
    G_t^{\mathcal{L}}(K{+}1)-G_t^{\mathcal{L}}(K)
    \;=\;
    \sigma_t\cdot\Delta\phi_K^{\mathcal{L}}
    + O_{K,d,L_t}(g_t^2 h),
    \qquad
    \Delta\phi_K^{\mathcal{L}}:=\phi_{K+1}^{\mathcal{L}}-\phi_K^{\mathcal{L}}.
    \label{eq:marginal-general}
  \end{equation}
  The dual-threshold stopping rule of \cref{alg:offline_online_budget}
  remains valid at the level of a proxy control law: stop when both the
  estimated marginal gain is small and the empirical within-timestep
  score variance is low, indicating low exhaustion-adjusted gain and a
  low sensitivity proxy.
\end{enumerate}
\end{corollary}

\begin{proof}
For part~(i): for each fixed allocation $\{K_t\}$, the map
$\boldsymbol{\sigma}\mapsto\sum_t\sigma_t\,\phi_{K_t}^{\mathcal{L}}$
is linear.
$V_{\mathcal{L}}^*$ is the pointwise maximum of finitely many
linear functions, hence convex.
Jensen's inequality gives~\eqref{eq:jensen-general}.
This argument is identical to \cref{sec:proof-online} and uses
\emph{no} property of $\phi_K^{\mathcal{L}}$ beyond non-negativity.

For part~(ii): the factorization~\eqref{eq:marginal-general} follows by
taking first differences in~\eqref{eq:general-gain-factorization}.
The dual-threshold logic of \cref{prop:online}(ii) relies solely on
the product structure of the marginal gain, the non-accelerating
property of $\phi_K^{\mathcal{L}}$, and the location-scale variance
proxy from \cref{thm:loc-scale}(ii); all three persist under
(A1)--(A4).
\end{proof}

\begin{remark}[Jensen gap magnitude and operator efficiency]
\label{rmk:jensen-curvature}
The Jensen gap~\eqref{eq:jensen-general} inherits the curvature of
$\phi_K^{\mathcal{L}}$.
Faster-growing $\phi_K$ (e.g., local perturbation with
$\phi_K^{(\mathrm{LP})}=c_d(\lambda)K$ in the linearized regime)
sharpens the switching between allocation regimes as
$\boldsymbol{\sigma}$ varies, enlarging the gap:
\emph{online adaptation is more valuable for more efficient
operators}.
For pure random search ($\phi_K=a_K\sim\sqrt{2\log K}$), the slow
growth yields a smaller gap and less benefit from online
scheduling.
\end{remark}

\begin{example}[Random search]
\label{ex:random-search}
The operator draws $K$ i.i.d.\ $\xi_k\sim\mathcal{N}(0,I)$
and keeps the best.
In the linearized regime, the projected scores
$\langle u,\xi_k\rangle$ are i.i.d.\ $\mathcal{N}(0,1)$, so
\begin{equation}
  \phi_K^{\mathrm{RS}} = a_K
  := \mathbb{E}\bigl[\max(Z_1,\ldots,Z_K)\bigr],
  \qquad Z_k\stackrel{\mathrm{iid}}{\sim}\mathcal{N}(0,1).
  \label{eq:phi-rs}
\end{equation}
\textbf{(A1):}~The running maximum is strictly increasing for $K\ge 2$
(a fresh draw has positive probability of exceeding the current max).
\textbf{(A2):}~$a_K$ is concave
(\citealt{david2003order}, Ch.~4).
\textbf{(A3):}~By \cref{thm:loc-scale}(iii),
$G_t(K)=\sigma_t\,a_K+O_{K,d,L_t}(g_t^2 h)$; the gain sequence
$a_K\sim\sqrt{2\log K}$ is $t$-independent.
\textbf{(A4):}~i.i.d.\ $\mathcal{N}(0,I)$ draws are
rotationally invariant.
\end{example}

\begin{example}[$\epsilon$-greedy search]
\label{ex:eps-greedy}
Starting from $\xi^{(0)}\sim\mathcal{N}(0,I)$, each iteration
draws $\xi_{\mathrm{new}}\sim\mathcal{N}(0,I)$ with probability
$\epsilon$ (explore), or sets
$\xi_{\mathrm{new}}=\xi^*+RU$ (exploit by perturbing the current best
$\xi^*$), where $U$ is uniform on the unit sphere $\mathbb{S}^{d-1}$ and
$R\sim\mathrm{Unif}[0,\lambda\sqrt d]$, and updates $\xi^*$ if the new
candidate is better.
In the linearized regime, the objective
$f(\xi)=\langle u,\xi\rangle$ is linear, and:
\begin{itemize}
  \item \emph{Exploration steps} ($\approx\epsilon K$ in expectation):
    each is an i.i.d.\ draw from $\mathcal{N}(0,1)$ in the projected
    coordinate $\langle u,\xi\rangle$, contributing a gain of
    $a_{\lceil\epsilon K\rceil}$.
  \item \emph{Exploitation steps} ($\approx(1{-}\epsilon)K$):
    each perturbs the current best in the projected direction by
    $R\langle u,U\rangle$.
    The accepted increment is $(R\langle u,U\rangle)_+$, whose
    expectation is the uniform-ball drift constant
    \begin{equation}
      c_d(\lambda)
      \;:=\;
      \mathbb{E}\bigl[(R\langle u,U\rangle)_+\bigr]
      \;=\;
      \frac{\lambda\sqrt d}{4\sqrt{\pi}}\,
      \frac{\Gamma(d/2)}{\Gamma((d+1)/2)}.
      \label{eq:uniform-ball-drift}
    \end{equation}
    Thus exploitation adds a linear-in-$K$ drift beyond the
    exploration component, although we retain only the robust lower
    bound below.
\end{itemize}
Therefore
\begin{equation}
  \phi_K^{(\epsilon)}
  \;\ge\; a_{\lceil\epsilon K\rceil},
  \label{eq:phi-epsgreedy-lb}
\end{equation}
with equality when exploitation steps add no improvement
(e.g., when $\lambda\to 0$).
In the other extreme, when $\epsilon=1$ the operator
reduces to random search and $\phi_K^{(1)}=a_K$.

\textbf{(A1):}~Each iteration has probability $\epsilon>0$ of being
an exploration step that draws a fresh
$Z_{\mathrm{new}}\sim\mathcal{N}(0,1)$ in the projected coordinate.
Because $\mathcal{N}(0,1)$ has unbounded support,
$P(Z_{\mathrm{new}}>\max_{0\le j\le K}\langle u,\xi^{(j)}\rangle)>0$,
so $\phi_{K+1}^{(\epsilon)}>\phi_K^{(\epsilon)}$ for all~$K\ge 0$.

\textbf{(A2):}~Let $M_K=\max_{0\le j\le K}\langle u,\xi^{(j)}\rangle$
denote the running best after $K$ iterations.
The marginal gain of the $(K{+}1)$-th iteration is
\[
  \Delta_K
  := \phi_{K+1}^{(\epsilon)}-\phi_K^{(\epsilon)}
  = \mathbb{E}\bigl[\max(0,\,\langle u,\xi^{(K+1)}\rangle - M_K)\bigr].
\]
With probability $\epsilon$ this iteration explores
(draws $Z_{\mathrm{new}}\sim\mathcal{N}(0,1)$ independent of $M_K$),
and with probability $1{-}\epsilon$ it exploits
(perturbs the incumbent by $R\langle u,U\rangle$).
In both cases, the candidate is accepted only when it exceeds~$M_K$.
For exploration, the candidate $Z_{\mathrm{new}}\sim\mathcal{N}(0,1)$
is independent of history, so
$\mathbb{E}[\max(0,Z_{\mathrm{new}}-M_K)]$ is non-increasing
in~$M_K$ (and hence in~$K$), because the tail probability
$P(Z_{\mathrm{new}}>m)$ decreases in~$m$.
For exploitation, the overshoot $(R\langle u,U\rangle)_+$
does not depend on~$M_K$ at all (the perturbation is added to
the incumbent and accepted only if its projected increment is positive),
yielding the constant marginal $c_d(\lambda)$ from
\eqref{eq:uniform-ball-drift}.
Therefore
$\Delta_K = \epsilon\cdot(\text{non-increasing in }K)
+ (1{-}\epsilon)\cdot(\text{constant})$,
which is non-increasing in~$K$,
so $\phi_K^{(\epsilon)}$ satisfies~(A2).

\textbf{(A3):}~In the linearized regime the projected problem
reduces to optimizing $\langle u,\xi\rangle$.
Both the exploration draw $\langle u,\xi_{\mathrm{new}}\rangle
\sim\mathcal{N}(0,1)$ and the exploitation perturbation
$R\langle u,U\rangle$ depend on the direction~$u$ only through a
one-dimensional projection, whose distribution is $u$-independent by
the isotropy of the Gaussian draw and of the sphere-uniform direction.
Hence $\phi_K^{(\epsilon)}$ does not depend on the timestep~$t$.

\textbf{(A4):}~Both primitives (Gaussian draws for explore,
sphere-uniform directions with scalar radii for exploit) are
rotationally invariant.
The accept-reject comparison uses
$\langle u,\xi\rangle$, which satisfies
$\langle Qu,Q\xi\rangle=\langle u,\xi\rangle$ for any
orthogonal~$Q$, so the full operator is rotationally equivariant.
\end{example}

\begin{example}[Local perturbation search ($\epsilon$-greedy with $\epsilon=0$)]
\label{ex:local-perturbation}
This operator is the pure-exploitation limit of the $\epsilon$-greedy
operator in \cref{ex:eps-greedy}: each iteration perturbs the
current best $\xi^*$ by a random direction
$U\sim\mathrm{Unif}(\mathbb{S}^{d-1})$ and a random radius
$R\sim\mathrm{Unif}[0,\lambda\sqrt d]$, setting
$\xi_{\mathrm{new}}=\xi^*+RU$, and updates $\xi^*$ if the
new candidate scores higher.
In the linearized regime, the projected coordinate of the
incumbent evolves by accept-reject: the perturbation
$R\langle u,U\rangle$ is
accepted when $\langle u,\xi_{\mathrm{new}}\rangle
>\langle u,\xi^*\rangle$, i.e., when $\langle u,U\rangle>0$.
Because the acceptance event is independent of the current
value of the incumbent, each iteration adds a constant expected
improvement
\begin{equation}
  \phi_{K+1}^{(\mathrm{LP})}-\phi_K^{(\mathrm{LP})}
  \;=\;
  c_d(\lambda),
  \label{eq:phi-lp-marginal}
\end{equation}
giving
\begin{equation}
  \phi_K^{(\mathrm{LP})}
  \;=\;
  c_d(\lambda)\,K.
  \label{eq:phi-lp-exact}
\end{equation}
Equivalently, if one parameterizes the perturbation by the radius cap
$\eta:=\lambda\sqrt d$, then
\[
  c_d(\eta)
  \;=\;
  \frac{\eta}{4\sqrt{\pi}}\,
  \frac{\Gamma(d/2)}{\Gamma((d+1)/2)}.
\]

\textbf{(A1):}~Each perturbation has probability
$P(\langle u,U\rangle>0)=\tfrac{1}{2}>0$ of improving the
incumbent, so $\phi_K$ is strictly increasing.

\textbf{(A2):}~In the linearized regime the marginal
gain~\eqref{eq:phi-lp-marginal} is constant (hence bounded
and non-increasing), satisfying~(A2).
Outside the linearized regime, curvature effects perturb this exact
constant-marginal picture; the abstract theory uses
\eqref{eq:phi-lp-exact} only as the leading-order gain sequence.

\textbf{(A3):}~The perturbation
$R\langle u,U\rangle$ is a one-dimensional projection of a
uniform-ball perturbation, whose distribution does not depend on~$u$
or the timestep~$t$.
Hence $\phi_K^{(\mathrm{LP})}$ is $t$-independent.

\textbf{(A4):}~The perturbation $RU$ is
rotationally invariant, and the accept-reject comparison depends
on $\langle u,\xi\rangle$ which satisfies
$\langle Qu,Q\xi\rangle=\langle u,\xi\rangle$ for any
orthogonal~$Q$.
\end{example}

\begin{proposition}[Asymptotic crossover scale between random search and local perturbation]
\label{prop:crossover}
In the linearized regime, random search achieves
$\phi_K^{\mathrm{RS}}\sim\sqrt{2\log K}$,
while local perturbation achieves
$\phi_K^{(\mathrm{LP})}= c_d(\lambda)K$
(\cref{eq:phi-lp-exact}).
Balancing the two growth laws yields a crossover scale
\begin{equation}
  K^*\;=\;\tilde{\Theta}\!\bigl(c_d(\lambda)^{-1}\bigr)
  \;=\;\tilde{\Theta}(\lambda^{-1})
  \label{eq:crossover}
\end{equation}
\noindent
where the final equivalence holds for fixed~$d$.
Thus budgets of order $K^*$ mark the regime in which the logarithmic
random-search gain and the linear local-perturbation gain are of the
same magnitude.
\end{proposition}

\begin{proof}
The crossover occurs at the unique solution of
$c_d(\lambda)K=\sqrt{2\log K}$.
Since $c_d(\lambda)K$ grows linearly and $\sqrt{2\log K}$
grows logarithmically, the balance point satisfies
$K=\sqrt{2\log K}/c_d(\lambda)$.
Substituting this relation once and
iterating once gives
$K^*=c_d(\lambda)^{-1}\sqrt{2\log(c_d(\lambda)^{-1})}\cdot(1+o(1))$
as $c_d(\lambda)\to 0$, which is
$\tilde{\Theta}(c_d(\lambda)^{-1})$.
For fixed dimension~$d$, the Gamma-ratio in
\eqref{eq:uniform-ball-drift} is constant, so
$c_d(\lambda)=\Theta(\lambda)$ and hence
$K^*=\tilde{\Theta}(\lambda^{-1})$.
\end{proof}

\begin{remark}[Full-problem caveat for crossover]
\label{rmk:crossover-full}
\cref{prop:crossover} compares only the linearized gain models.
In the full nonlinear objective, the remainder
$R_t(\xi)=O(g_t^2 h\|\xi\|^2)$ couples performance to the
operator-dependent iterate norms, so any large-budget ordering beyond
the crossover scale becomes problem dependent and is best treated as an
implementation-level heuristic rather than a theorem.
\end{remark}

\begin{remark}[Operator selection as part of scheduling]
\label{rmk:operator-selection}
\cref{prop:crossover} implies that the optimal operator choice
depends on the per-step budget $K_t$: as $K_t$ approaches the crossover
scale, the linear-growth local perturbation operator becomes relatively
more attractive, while smaller budgets remain well served by pure random
search.
Extending the scheduling decision to include operator selection is a
natural next step; the MDP formulation of \cref{sec:mdp} already
accommodates this as a richer action space.
\end{remark}

\bibliographystyle{plainnat}
\bibliography{main}

\end{document}